\newcommand{\titulus}{\nomenFesti{S. Andreæ Dũng Lac, Presbyteri \& Sociorum Martyrum.}
\dies{Die 24. Novembris.}}
\newcommand{\oratio}{\pars{Oratio.}

\noindent Deus, omnis paternitátis fons et orígo, qui beátos mártyres Andréam et sócios eius Cruci Fílii tui usque ad sánguinis effusiónem fidéles effecísti, eórum intercessióne concéde, ut amórem tuum inter fratres propagántes fílii tui nominári et esse valeámus.

\pars{Pro pace in universo mundo.} \scriptura{Sir. 50, 25; 2 Esdr. 4, 20; \textbf{H416}}

\vspace{-4mm}

\antiphona{II D}{temporalia/ant-dapacemdomine.gtex}

\vfill

\noindent Deus, a quo sancta desidéria, recta consília et iusta sunt ópera: da servis tuis illam, quam mundus dare non potest, pacem; ut et corda nostra mandátis tuis dédita, et hóstium subláta formídine, témpora sint tua protectióne tranquílla.

\noindent Per Dóminum nostrum Iesum Christum, Fílium tuum, qui tecum vivit et regnat in unitáte Spíritus Sancti, Deus, per ómnia sǽcula sæculórum.

\noindent \Rbardot{} Amen.}
\newcommand{\invitatorium}{\pars{Invitatorium.}

\vspace{-4mm}

\antiphona{E}{temporalia/inv-regemmartyrumsimplex.gtex}}
\newcommand{\hymnusmatutinum}{\pars{Hymnus}

\cuminitiali{IV}{temporalia/hym-RexGloriose.gtex}}
\newcommand{\lectioi}{\pars{Lectio I.} \scriptura{2 Petr. 1, 12-21}

\noindent De Epístola secúnda beáti Petri apóstoli.

\noindent Caríssimi: Incípiam vos semper commonére de his, et quidem sciéntes et confirmátos in præsénti veritáte. Iustum autem árbitror, quámdiu sum in hoc tabernáculo, suscitáre vos in commonitióne, certus quod velox est deposítio tabernáculi mei, secúndum quod et Dóminus noster Iesus Christus significávit mihi; dabo autem óperam et frequénter habére vos post óbitum meum, ut horum memóriam faciátis.

\noindent Non enim captiósas fábulas secúti notam fécimus vobis Dómini nostri Iesu Christi virtútem et advéntum, sed speculatóres facti illíus magnitúdinis. Accípiens enim a Deo Patre honórem et glóriam, voce proláta ad eum huiuscémodi a magnífica glória: «Fílius meus, diléctus meus hic est, in quo ego mihi complácui»; et hanc vocem nos audívimus de cælo prolátam, cum essémus cum ipso in monte sancto. Et habémus firmiórem prophéticum sermónem, cui bene fácitis attendéntes quasi lucérnæ lucénti in caliginóso loco, donec dies illucéscat, et lúcifer oriátur in córdibus vestris, hoc primum intellegéntes quod omnis prophetía Scriptúræ própria interpretatióne non fit; non enim voluntáte humána proláta est prophetía aliquándo, sed a Spíritu Sancto ducti locúti sunt a Deo hómines.}
\newcommand{\responsoriumi}{\pars{Responsorium 1.} \scriptura{\Rbardot{} Ier. 31, 11.12 \Vbardot{} Ps. 4, 8; \textbf{H418}}

\vspace{-5mm}

\responsorium{III}{temporalia/resp-redemitdominus-CROCHU.gtex}{}}
\newcommand{\lectioii}{\pars{Lectio II.} \scriptura{A. Launay, Le clergé tonkinois et ses prêtres martyrs, MEP, Paris 1925, pp. 80-83}

\noindent E lítteris sancti Pauli Le-Bao-Tinh {\color{gray} alúmnis Seminárii Ke-Vinh anno 1843 expedítis.}

\noindent Ego, Paulus, pro nómine Christi vinctus, tribulatiónes meas vobis reférre volo quibus cotídie immérsus sum, ita ut, amóre erga Deum accénsi laudes mecum Deo præbeátis, \emph{quóniam in ætérnum misericórdia eius.} Hic carcer vere imágo est inférni ætérni: ad supplícia crudélia omnis géneris, ut sunt cómpedes, caténæ férreæ et víncula, addúntur ódium, vindíctæ, calúmniæ, verba indecéntia, querélæ, actus mali, iuraménta iniústa, maledictiónes et tandem angústiæ et tristítia. Deus autem qui olim liberávit tres púeros de camíno ignis, mihi semper adest meque ab istis tribulatiónibus liberávit et eas in dulcédinem convértit, \emph{quóniam in ætérnum misericórdia eius.}

\noindent In médio autem horum tormentórum, quæ álios contérrere solent, grátia Dei, gáudio replétus sum et lætítia, quia non solus sed cum Christo sum.

\noindent Ipse magíster noster totum pondus crucis sústinet, mihi mínimam tantum et últimam partem impónens. Certáminis mei non solum spectátor, sed ipse est bellátor et victor totiúsque agónis consummátor. Proptérea super caput eius pósita est coróna victóriæ, cuius glóriam partícipant étiam eius membra.

\noindent Quómodo autem sustíneam spectáculum istud, videns cotídie imperatóres, mandarínos eorúmque satéllites blasphemántes nomen sanctum tuum, Dómine, \emph{qui sedes super Chérubim}et Séraphim? Ecce, crux tua a pédibus paganórum conculcáta est! Ubi est glória tua? Videns hæc ómnia, malo, amóre tui succénsus, abscíssis membris, mori in testimónium amóris tui.

\noindent Osténde, Dómine, poténtiam tuam, salva me et sústine me, ut \emph{virtus in infirmitáte} mea ostendátur et glorificétur coram géntibus, ne, cum vacíllem forte in via, inimíci tui in supérbia sua caput possint eleváre.}
\newcommand{\responsoriumii}{\pars{Responsorium 2.} \scriptura{\Rbardot{} Ps. 149, 2 \Vbardot{} Ps. 67, 4; \textbf{H369}}

\vspace{-5mm}

\responsorium{VII}{temporalia/resp-exsultabuntsancti-CROCHU.gtex}{}}
\newcommand{\lectioiii}{\pars{Lectio III.}

\noindent Fratres caríssimi, audiéntes hæc ómnia, grátias agátis immortáles in lætítia Deo, a quo bona cuncta procédunt, benedícite Dómino mecum, \emph{quóniam in ætérnum misericórdia eius!} Magníficet ánima mea Dóminum et exsúltet spíritus meus in Deo meo, quóniam respéxit humilitátem fámuli sui et ex hoc beátum me dicent omnes generatiónes futúræ: \emph{quóniam in ætérnum misericórdia eius.}

\noindent \emph{Laudáte Dóminum omnes gentes, collaudáte eum omnes pópuli,} quóniam quæ \emph{infírma sunt mundi, elégit Deus ut confúndat fórtia et ignobília mundi et contemptibília elégit Deus} ut nobília confúndat. Per os meum atque intelléctum meum confúdit philósophos qui discípuli sunt sapiéntum huius mundi, \emph{quóniam in ætérnum misericórdia eius.}

\noindent Scribo vobis hæc ómnia, ut uniántur fides vestra et mea. In médio huius tempestátis áncoram iácio usque ad thronum Dei; spem vivam, quæ est in corde meo.

\noindent Vos autem, fratres caríssimi, \emph{sic cúrrite ut corónam comprehendátis,} indúite \emph{lorícam fídei} et \emph{arma} Christi súmite \emph{a dextris et a sinístris,} sicut dócuit sanctus Paulus, patrónus meus. \emph{Bonum vobis est, unóculos vel débiles in vitam intráre,} quam ómnia membra habéntes foris mitti.

\noindent Succúrrite mihi précibus vestris, ut secúndum legem certáre, et quidem \emph{bonum certámen certáre} et usque in finem certáre váleam, cursum meum felíciter consummatúrus; si in hac vita non iam nos vidébimus, in futúro tamen sǽculo hæc felícitas nostra erit, quando astántes ad thronum Agni immaculáti, unánimes laudes eius cantábimus exsultántes in gáudio victóriæ in perpétuum. Amen.}
\newcommand{\responsoriumiii}{\pars{Responsorium 3.} \scriptura{Cf. Ap. 14, 4; \textbf{H68}}

\vspace{-5mm}

\responsorium{III}{temporalia/resp-coronaviteos-CROCHU-cumdox.gtex}{}}
\newcommand{\hymnuslaudes}{\pars{Hymnus}

\cuminitiali{VIII}{temporalia/hym-AEternaChristi.gtex}}
\newcommand{\lectiobrevis}{\pars{Lectio Brevis.} \scriptura{2 Cor. 1, 3-5}

\noindent Benedíctus Deus et Pater Dómini nostri Iesu Christi, Pater misericordiárum et Deus totíus consolatiónis, qui consolátur nos in omni tribulatióne nostra, ut possímus et ipsi consolári eos, qui in omni pressúra sunt, per exhortatiónem, qua exhortámur et ipsi a Deo; quóniam, sicut abúndant passiónes Christi in nobis, ita per Christum abúndat et consolátio nostra.}
\newcommand{\responsoriumbreve}{\pars{Responsorium breve.} \scriptura{Sap. 5, 16}

\antiphona{VI}{temporalia/resp-iustiautem.gtex}}
\newcommand{\benedictus}{\pars{Canticum Zachariæ.} \scriptura{Ap. 21, 4}

\vspace{-4mm}

\antiphona{I f}{temporalia/ant-abstergetdeus.gtex}

\vspace{-2mm}

\scriptura{Lc. 1, 68-79}

\vspace{-2mm}

\cantusSineNeumas
\initiumpsalmi{temporalia/benedictus-initium-i-f-auto.gtex}

%\vspace{-1.5mm}

\input{temporalia/benedictus-i-f.tex} \Abardot{}}
\newcommand{\preces}{\noindent Fratres, Salvatórem nostrum, testem fidélem, per mártyres interféctos propter verbum Dei, celebrémus,~\grestar{} clamántes:

\Rbardot{} Redemísti nos Deo in sánguine tuo.

\noindent Per mártyres tuos, qui líbere mortem in testimónium fídei sunt ampléxi,~\grestar{} da nobis, Dómine, veram spíritus libertátem.

\Rbardot{} Redemísti nos Deo in sánguine tuo.

\noindent Per mártyres tuos, qui fidem usque ad sánguinem sunt conféssi,~\grestar{} da nobis, Dómine, puritátem fideíque constántiam.

\Rbardot{} Redemísti nos Deo in sánguine tuo.

\noindent Per mártyres tuos, qui, sustinéntes crucem, tua vestígia sunt secúti,~\grestar{} da nobis, Dómine, ærúmnas vitæ fórtiter sustinére.

\Rbardot{} Redemísti nos Deo in sánguine tuo.

\noindent Per mártyres tuos, qui stolas suas lavérunt in sánguine Agni,~\grestar{} da nobis, Dómine, omnes insídias carnis mundíque devíncere.

\Rbardot{} Redemísti nos Deo in sánguine tuo.}
\newcommand{\benedicamuslaudes}{\cuminitiali{}{temporalia/benedicamus-memoria-laudes.gtex}}
\newcommand{\hebdomada}{infra Hebdom. XXXIV per Annum.}
\newcommand{\hiemalis}{Hiemalis}
\newcommand{\matub}{Matutinum Hebdomadae B}
\newcommand{\matubd}{Matutinum Hebdomadae B vel D}
\newcommand{\laudb}{Laudes Hebdomadae B}
\newcommand{\laudbd}{Laudes Hebdomadae B vel D}

% LuaLaTeX

\documentclass[a4paper, twoside, 12pt]{article}
\usepackage[latin]{babel}
%\usepackage[landscape, left=3cm, right=1.5cm, top=2cm, bottom=1cm]{geometry} % okraje stranky
%\usepackage[landscape, a4paper, mag=1166, truedimen, left=2cm, right=1.5cm, top=1.6cm, bottom=0.95cm]{geometry} % okraje stranky
\usepackage[landscape, a4paper, mag=1400, truedimen, left=0.5cm, right=0.5cm, top=0.5cm, bottom=0.5cm]{geometry} % okraje stranky

\usepackage{fontspec}
\setmainfont[FeatureFile={junicode.fea}, Ligatures={Common, TeX}, RawFeature=+fixi]{Junicode}
%\setmainfont{Junicode}

% shortcut for Junicode without ligatures (for the Czech texts)
\newfontfamily\nlfont[FeatureFile={junicode.fea}, Ligatures={Common, TeX}, RawFeature=+fixi]{Junicode}

\usepackage{multicol}
\usepackage{color}
\usepackage{lettrine}
\usepackage{fancyhdr}

% usual packages loading:
\usepackage{luatextra}
\usepackage{graphicx} % support the \includegraphics command and options
\usepackage{gregoriotex} % for gregorio score inclusion
\usepackage{gregoriosyms}
\usepackage{wrapfig} % figures wrapped by the text
\usepackage{parcolumns}
\usepackage[contents={},opacity=1,scale=1,color=black]{background}
\usepackage{tikzpagenodes}
\usepackage{calc}
\usepackage{longtable}
\usetikzlibrary{calc}

\setlength{\headheight}{14.5pt}

% Commands used to produce a typical "Conventus" booklet

\newenvironment{titulusOfficii}{\begin{center}}{\end{center}}
\newcommand{\dies}[1]{#1

}
\newcommand{\nomenFesti}[1]{\textbf{\Large #1}

}
\newcommand{\celebratio}[1]{#1

}

%%%% Prejmenovat na latinske
\newcounter{Thumb}
\setcounter{Thumb}{0}
\newcommand{\thumbTitulus}{}

\newif\ifSideThumbsEnabled

\newlength\ThumbSize
\setlength\ThumbSize{2cm}

\SideThumbsEnabledfalse
\AddEverypageHook{
\ifSideThumbsEnabled
\ifodd\value{page}
\backgroundsetup{
angle=90,
position={current page.east|-current page text area.north east},
vshift=12pt,
hshift=-\value{Thumb}*\ThumbSize+\ThumbSize+\ThumbSize/4,
contents={\tikz\node[fill=gray!30,anchor=west,text width=\ThumbSize,
 align=center,text height=14pt,text depth=30pt]
{\small\thumbTitulus};
}
}
\else
\backgroundsetup{
angle=90,
position={current page.west|-current page text area.north west},
vshift=-16pt,
hshift=-\value{Thumb}*\ThumbSize+\ThumbSize+\ThumbSize/4,
contents={\tikz\node[fill=gray!30,anchor=west,text width=\ThumbSize,
 align=center,text height=28pt,text depth=6pt]
{\small\thumbTitulus};
}
}
\fi
\BgMaterial
\else\relax\fi}
\newcommand\sideThumbs[1]{\SideThumbsEnabledtrue%
\renewcommand{\thumbTitulus}{#1}%
\stepcounter{Thumb}
}
\newcommand\RemoveSideThumbs{\SideThumbsEnabledfalse}

\newcommand{\hora}[1]{%
\vspace{0.5cm}{\large \textbf{#1}}

\fancyhead[LE]{\thepage\ / #1}
\fancyhead[RO]{#1 / \thepage}
\addcontentsline{toc}{subsection}{#1}
}

% In Junicode, the default R/ with .35em has the bar too close to the
% letter.
\grelatexsimpledefbarredsymbol{A}{.30em}{.25em}{.30em}{.27em}%
\grelatexsimpledefbarredsymbol{R}{.41em}{.45em}{.45em}{.48em}%
\grelatexsimpledefbarredsymbol{V}{.10em}{.10em}{.10em}{.13em}%

% Emit R/, V/, A/ and + in red
\let\oldrbar\Rbar
\let\oldvbar\Vbar
\let\oldabar\Abar
\let\oldgrecross\grecross
\let\oldgrealtcross\grealtcross
\let\oldGreStar\GreStar
\let\oldGreDagger\GreDagger
\let\olddag\dag
\let\oldddag\ddag
\renewcommand{\Rbar}{{\color{red}\oldrbar}}
\renewcommand{\Vbar}{{\color{red}\oldvbar}}
\renewcommand{\Abar}{{\color{red}\oldabar}}
\renewcommand{\grecross}{{\color{red}\oldgrecross}}
\renewcommand{\grealtcross}{{\color{red}\oldgrealtcross}}
\newcommand{\grestar}{{\color{red}\oldGreStar}}
\renewcommand{\GreStar}{{\color{red}\oldGreStar}}
\newcommand{\gredagger}{{\color{red}\oldGreDagger}}
\renewcommand{\GreDagger}{{\color{red}\oldGreDagger}}
\renewcommand{\dag}{{\color{red}\olddag}}
\renewcommand{\ddag}{{\color{red}\oldddag}}
\newcommand{\Rbardot}{{\color{red}\oldrbar.~}}
\newcommand{\Vbardot}{{\color{red}\oldvbar.~}}
\newcommand{\Abardot}{{\color{red}\oldabar.~}}
\newcommand{\Psdot}{{\color{red}Ps.~}}

\let\quasiHora\hora

% larger unit than a hora
\newcommand{\divisio}[1]{%
\begin{center}
{\Large \textsc{#1}}
\end{center}
\fancyhead[CO,CE]{#1}
\addcontentsline{toc}{section}{#1}
}

% a part of a hora, larger than pars
\newcommand{\subhora}[1]{%
\begin{center}
{\large \textit{#1}}
\end{center}
%\fancyhead[CO,CE]{#1}
\addcontentsline{toc}{subsubsection}{#1}
}

% rubricated inline text
\newcommand{\rubricatum}[1]{\textit{#1}}

% standalone rubric
\newcommand{\rubrica}[1]{\vspace{3mm}\rubricatum{#1}}

\newcommand{\notitia}[1]{\textcolor{red}{#1}}

\newcommand{\scriptura}[1]{\hfill \small\textit{#1}}

\newcommand{\translatioCantus}[1]{\vspace{1mm}%
{\noindent\footnotesize \nlfont{#1}}}

\newenvironment{translatioMulticol}[1]{%
  \begin{multicols}{#1}
    \footnotesize
    \setlength{\parindent}{0pt}
    \setlength{\columnseprule}{0pt}
}{\end{multicols}}

% pruznejsi varianta nasledujiciho - umoznuje nastavit sirku sloupce
% s prekladem
\newcommand{\psalmusEtTranslatioB}[3]{%
  \vspace{0.5cm}
  \begin{parcolumns}[colwidths={2=#3}, nofirstindent=true]{2}
    \colchunk{%
      \input{#1}
    }

    \colchunk{%
      \vspace{-0.5cm}
      {\footnotesize \nlfont
        \input{#2}
      }
    }
  \end{parcolumns}
}

\newcommand{\psalmusEtTranslatio}[2]{%
  \psalmusEtTranslatioB{#1}{#2}{8.5cm}
}

\newcommand{\psalmusTr}[1]{{\footnotesize\nlfont\hspace*{-4pt} #1}}
\newcommand{\psalmusEtTranslatioT}[2]{{%
  \setlength{\LTleft}{-0.25cm}
  \setlength{\LTright}{\fill}
  \setlength{\tabcolsep}{0.35cm}
  \begin{psalmus}
      \noindent
      \begin{longtable}{p{\textwidth-#2-1cm}p{#2}}
        \input{#1}
      \end{longtable}
  \end{psalmus}
}}
\newcommand{\psalmusEtTranslatioTS}[2]{{%
  \setlength{\LTleft}{-0.25cm}
  \setlength{\LTright}{\fill}
  \setlength{\tabcolsep}{0.25cm}
  \begin{psalmus}
      \noindent
      \begin{longtable}{p{\textwidth-#2-0.5cm}p{#2}}
        \input{#1}
      \end{longtable}
  \end{psalmus}
}}

\newcommand{\canticumMagnificatEtTranslatio}[1]{%
  \psalmusEtTranslatioB{#1}{temporalia/extra-adventum-vespers/magnificat-boh.tex}{12cm}
}
\newcommand{\canticumBenedictusEtTranslatio}[1]{%
  \psalmusEtTranslatioB{#1}{temporalia/extra-adventum-laudes/benedictus-boh.tex}{10.5cm}
}

\newcommand{\textusEtTranslatio}[3]{%
  \begin{parcolumns}[colwidths={2=#3}, nofirstindent=true]{2}
    \colchunk{%
      #1
    }

    \colchunk{%
      \vspace{-0.5cm}
      {\footnotesize \nlfont
        #2
      }
    }
  \end{parcolumns}
}

% volne misto nad antifonami, kam si zpevaci dokresli neumy
\newcommand{\hicSuntNeumae}{\vspace{0.5cm}}

% prepinani mista mezi notovymi osnovami: pro neumovane a neneumovane zpevy
\newcommand{\cantusCumNeumis}{%
  \grechangestaffsize{14}
%  \global\advance\grespaceabovelines by 5mm%
%  \global\advance\grespaceabovelines by 2mm%
}
\newcommand{\cantusSineNeumas}{%
  \grechangestaffsize{14}
%  \global\advance\grespaceabovelines by -5mm%
  \grechangedim{spaceabovelines}{0cm}{scalable}%
}

\grechangestyle{abovelinestext}{\color{red}\it}

% znaky k umisteni nad inicialu zpevu
\newcommand{\superInitialam}[1]{\greannotation{\footnotesize {\textbf{#1}}}}

% pars officii, i.e. "oratio", ...
\newcommand{\pars}[1]{\textbf{#1}}

\newenvironment{psalmus}{%
  \setlength{\parindent}{0pt}
  \setlength{\parskip}{5pt}
}{%
  \setlength{\parindent}{10pt}
  \setlength{\parskip}{10pt}
}

%%%% Prejmenovat na latinske:
\newcommand{\nadpisZalmu}[1]{%
  \hspace{2cm}\textbf{#1}\vspace{2mm}%
  \nopagebreak%

}

% mode, score, translation
\newcommand{\antiphona}[3]{%
  \hicSuntNeumae%
  \superInitialam{#1}%
  \gresetinitiallines{1}%
  \gregorioscore{#2}%

#3
}

% a semantic alias of the previous
\newcommand{\responsorium}[3]{\antiphona{#1}{#2}{#3}}

\newcommand{\initiumpsalmi}[1]{%
  \gresetinitiallines{0}%
  \gregorioscore{#1}%
}

\newcommand{\cuminitiali}[2]{%
  \superInitialam{#1}%
  \gresetinitiallines{1}%
  \gregorioscore{#2}%
}

\newcommand{\sineinitiali}[1]{%
  \gresetinitiallines{0}%
  \gregorioscore{#1}%
}
 % Often used macros

\newcommand{\annusEditionis}{2022}

%%%% Vicekrat opakovane kousky

\newcommand{\anteOrationem}{
  \rubrica{Ante Orationem, cantatur a Superiore:}

  \pars{Supplicatio Litaniæ.}

  \cuminitiali{}{temporalia/supplicatiolitaniae.gtex}

  \pars{Oratio Dominica.}

  \cuminitiali{}{temporalia/oratiodominica.gtex}

  \rubrica{Deinde dicitur ab Hebdomadario:}

  \cuminitiali{}{temporalia/dominusvobiscum-solemnis.gtex}

  \rubrica{In choro monialium loco Dominus vobiscum dicitur:}

  \sineinitiali{temporalia/domineexaudi.gtex}
}

\setlength{\columnsep}{30pt} % prostor mezi sloupci

%%%%%%%%%%%%%%%%%%%%%%%%%%%%%%%%%%%%%%%%%%%%%%%%%%%%%%%%%%%%%%%%%%%%%%%%%%%%%%%%%%%%%%%%%%%%%%%%%%%%%%%%%%%%%
\begin{document}

% Here we set the space around the initial.
% Please report to http://home.gna.org/gregorio/gregoriotex/details for more details and options
\grechangedim{afterinitialshift}{2.2mm}{scalable}
\grechangedim{beforeinitialshift}{2.2mm}{scalable}
\grechangedim{interwordspacetext}{0.22 cm plus 0.15 cm minus 0.05 cm}{scalable}%
\grechangedim{annotationraise}{-0.2cm}{scalable}

% Here we set the initial font. Change 38 if you want a bigger initial.
% Emit the initials in red.
\grechangestyle{initial}{\color{red}\fontsize{38}{38}\selectfont}

\pagestyle{empty}

%%%% Titulni stranka
\begin{titulusOfficii}
\ifx\titulus\undefined
\nomenFesti{Feria III \hebdomada{}}
\else
\titulus
\fi
\end{titulusOfficii}

\vfill

\begin{center}
%Ad usum et secundum consuetudines chori \guillemotright{}Conventus Choralis\guillemotleft.

%Editio Sancti Wolfgangi \annusEditionis
\end{center}

\scriptura{}

\pars{}

\pagebreak

\renewcommand{\headrulewidth}{0pt} % no horiz. rule at the header
\fancyhf{}
\pagestyle{fancy}

\cantusSineNeumas

\ifx\oratio\undefined
\ifx\lauda\undefined
\else
\newcommand{\oratio}{\pars{Oratio.}

\noindent Matutína súpplicum vota, Dómine, propítius intuére et occúlta cordis nostri remédio tuæ clarífica pietátis, ut desidéria tenebrósa non téneant, quos lux cæléstis grátiæ reparávit. 

\noindent Per Dóminum nostrum Iesum Christum, Fílium tuum, qui tecum vivit et regnat in unitáte Spíritus Sancti, Deus, per ómnia sǽcula sæculórum.

\noindent \Rbardot{} Amen.}
\fi
\ifx\laudb\undefined
\else
\newcommand{\oratio}{\pars{Oratio.}

\noindent Dómine Iesu Christe, lux vera, qui omnes hómines illúminas ad salútem, nobis, quǽsumus, concéde virtútem, ut ante te vias pacis et iustítiæ præparémus.

\noindent Qui vivis et regnas cum Deo Patre in unitáte Spíritus Sancti, Deus, per ómnia sǽcula sæculórum.

\noindent \Rbardot{} Amen.}
\fi
\ifx\laudc\undefined
\else
\newcommand{\oratio}{\pars{Oratio.}

\noindent Deus omnípotens, cuius est bona rerum et speciósa creátio, da nobis hunc diem in nómine tuo lætánter incípere et in operósa tui fratrúmque dilectióne complére.

\noindent Per Dóminum nostrum Iesum Christum, Fílium tuum, qui tecum vivit et regnat in unitáte Spíritus Sancti, Deus, per ómnia sǽcula sæculórum.

\noindent \Rbardot{} Amen.}
\fi
\ifx\laudd\undefined
\else
\newcommand{\oratio}{\pars{Oratio.}

\noindent Auge in nobis, quǽsumus, Dómine, fidem tuam, ut in ore nostro perfécta laus tua cæléstes iúgiter áfferat fructus.

\noindent Per Dóminum nostrum Iesum Christum, Fílium tuum, qui tecum vivit et regnat in unitáte Spíritus Sancti, Deus, per ómnia sǽcula sæculórum.

\noindent \Rbardot{} Amen.}
\fi
\fi

\hora{Ad Matutinum.} %%%%%%%%%%%%%%%%%%%%%%%%%%%%%%%%%%%%%%%%%%%%%%%%%%%%%

\vspace{2mm}

\cuminitiali{}{temporalia/dominelabiamea.gtex}

\vfill
%\pagebreak

\vspace{2mm}

\ifx\invitatorium\undefined
%\ifx\matuac\undefined
%\else
%\pars{Invitatorium.} \scriptura{Ps. 94, 1; Psalmus 94; \textbf{H451}}
%
%\vspace{-6mm}
%
%\antiphona{VI}{temporalia/inv-jubilemusdeo.gtex}
%\fi
%\ifx\matubd\undefined
%\else
\pars{Invitatorium.} \scriptura{Cantor; Psalmus 94; \textbf{H449}}

\vspace{-6mm}

\antiphona{E}{temporalia/inv-regemmagnum.gtex}
%\fi
\else
\invitatorium
\fi

\vfill
\pagebreak

\ifx\hymnusmatutinum\undefined
\ifx\hiemalis\undefined
\ifx\matuac\undefined
\else
\pars{Hymnus.}

\antiphona{IV}{temporalia/hym-OSacrosanctaTrinitas-kempten.gtex}
\fi
\ifx\matub\undefined
\else
\pars{Hymnus.}

\antiphona{VIII}{temporalia/hym-AdPrecesNostras-kempten.gtex}
\fi
\ifx\matud\undefined
\else
\pars{Hymnus.}

\antiphona{I}{temporalia/hym-AdPrecesNostras.gtex}
\fi
\else
\ifx\matua\undefined
\else
\pars{Hymnus}

\cuminitiali{IV}{temporalia/hym-ConsorsPaterni.gtex}
\fi
\ifx\matuc\undefined
\else
\pars{Hymnus}

\cuminitiali{IV}{temporalia/hym-ConsorsPaterni-MMMA.gtex}
\fi
\ifx\matubd\undefined
\else
\pars{Hymnus.} \scriptura{Gregorius Magnus (\olddag{} 604)}

{
\grechangedim{interwordspacetext}{0.10 cm plus 0.15 cm minus 0.05 cm}{scalable}%
\antiphona{I}{temporalia/hym-NocteSurgentes.gtex}
\grechangedim{interwordspacetext}{0.22 cm plus 0.15 cm minus 0.05 cm}{scalable}%
}
\fi
\fi
\else
\hymnusmatutinum
\fi

\vspace{-3mm}

\vfill
\pagebreak

\ifx\matutinum\undefined
\ifx\matua\undefined
\else
% MAT A
\pars{Psalmus 1.} \scriptura{Ps. 9, 22}

\vspace{-4mm}

\antiphona{II D}{temporalia/ant-utquiddomine.gtex}

%\vspace{-2mm}

\scriptura{Ps. 9, 22-32}

%\vspace{-2mm}

\initiumpsalmi{temporalia/ps9xxii_xxxii-initium-ii-D-auto.gtex}

\input{temporalia/ps9xxii_xxxii-ii-D.tex} \Abardot{}

\vfill
\pagebreak

\pars{Psalmus 2.} \scriptura{Ex. 15, 18}

\vspace{-4mm}

\antiphona{IV* e}{temporalia/ant-inaeternum.gtex}

%\vspace{-2mm}

\scriptura{Ps. 9, 33-39}

%\vspace{-2mm}

\initiumpsalmi{temporalia/ps9xxxiii_xxxix-initium-iv_-e-auto.gtex}

\input{temporalia/ps9xxxiii_xxxix-iv_-e.tex} \Abardot{}

\vfill
\pagebreak

\pars{Psalmus 3.} \scriptura{Ps. 11, 8}

\vspace{-4mm}

\antiphona{VIII G}{temporalia/ant-tudomine.gtex}

%\vspace{-2mm}

\scriptura{Ps. 11}

%\vspace{-2mm}

\initiumpsalmi{temporalia/ps11-initium-viii-G-auto.gtex}

\input{temporalia/ps11-viii-G.tex} \Abardot{}

\vfill
\pagebreak
\fi
\ifx\matub\undefined
\else
% MAT B
\pars{Psalmus 1.} \scriptura{Ps. 36, 5; \textbf{H93}}

\vspace{-4mm}

\antiphona{VI F}{temporalia/ant-reveladomino.gtex}

%\vspace{-2mm}

\scriptura{Ps. 36, 1-11}

%\vspace{-2mm}

\initiumpsalmi{temporalia/ps36i_xi-initium-vi-F-auto.gtex}

\input{temporalia/ps36i_xi-vi-F.tex} \Abardot{}

\vfill
\pagebreak

\pars{Psalmus 2.}

\vspace{-4mm}

\antiphona{II D}{temporalia/ant-iuniorfui.gtex}

\vspace{-2mm}

\scriptura{Ps. 36, 12-29}

\vspace{-2mm}

\initiumpsalmi{temporalia/ps36xii_xxix-initium-ii-D-auto.gtex}

\input{temporalia/ps36xii_xxix-ii-D.tex}

\vfill

\antiphona{}{temporalia/ant-iuniorfui.gtex}

\vfill
\pagebreak

\pars{Psalmus 3.} \scriptura{Ps. 36, 3}

\vspace{-4mm}

\antiphona{VI F}{temporalia/ant-speraindomino.gtex}

%\vspace{-2mm}

\scriptura{Ps. 36, 30-40}

%\vspace{-2mm}

\initiumpsalmi{temporalia/ps36iii-initium-vi-F-auto.gtex}

\input{temporalia/ps36iii-vi-F.tex} \Abardot{}

\vfill
\pagebreak
\fi
\ifx\matuc\undefined
\else
% MAT C
\pars{Psalmus 1.} \scriptura{Ps. 67, 2}

\vspace{-4mm}

\antiphona{VII a}{temporalia/ant-exsurgatdeus.gtex}

%\vspace{-2mm}

\scriptura{Ps. 67, 2-11}

\initiumpsalmi{temporalia/ps67i-initium-vii-a-auto.gtex}

\input{temporalia/ps67i-vii-a.tex} \Abardot{}

\vfill
\pagebreak

\pars{Psalmus 2.}

\vspace{-4mm}

\antiphona{I f}{temporalia/ant-deusnosterdeussalvos.gtex}

%\vspace{-2mm}

\scriptura{Ps. 67, 12-24}

%\vspace{-2mm}

\initiumpsalmi{temporalia/ps67ii-initium-i-f-auto.gtex}

\input{temporalia/ps67ii-i-f.tex} \Abardot{}

\vfill
\pagebreak

\pars{Psalmus 3.} \scriptura{Ps. 67, 27; \textbf{H96}}

\vspace{-4mm}

\antiphona{D}{temporalia/ant-inecclesiis.gtex}

%\vspace{-2mm}

\scriptura{Ps. 67, 25-36}

\initiumpsalmi{temporalia/ps67iii-initium-d-g2-auto.gtex}

\input{temporalia/ps67iii-d-g2.tex} \Abardot{}

\vfill
\pagebreak
\fi
\ifx\matud\undefined
\else
% MAT D
\pars{Psalmus 1.} \scriptura{Ps. 101, 2; \textbf{H99}}

\vspace{-4mm}

\antiphona{I a}{temporalia/ant-clamormeus1.gtex}

%\vspace{-2mm}

\scriptura{Ps. 101, 2-12}

%\vspace{-2mm}

\initiumpsalmi{temporalia/ps101ii_xii-initium-i-a.gtex}

\input{temporalia/ps101ii_xii-i-a.tex} \Abardot{}

\vfill
\pagebreak

\pars{Psalmus 2.}

\vspace{-4mm}

\antiphona{VIII G}{temporalia/ant-respicehumilitatem.gtex}

\vspace{-2mm}

\scriptura{Ps. 101, 13-23}

\vspace{-2mm}

\initiumpsalmi{temporalia/ps101xiii_xxiii-initium-viii-G-auto.gtex}

\input{temporalia/ps101xiii_xxiii-viii-G.tex} \Abardot{}

\vfill
\pagebreak

\pars{Psalmus 3.} \scriptura{Ps. 139, 14; \textbf{H371}}

\vspace{-4mm}

\antiphona{I g}{temporalia/ant-iusticonfitebuntur1.gtex}

%\vspace{-2mm}

\scriptura{Ps. 101, 24-29}

%\vspace{-2mm}

\initiumpsalmi{temporalia/ps101iii-initium-i-g-auto.gtex}

\input{temporalia/ps101iii-i-g.tex} \Abardot{}
\vfill
\pagebreak
\fi
\else
\matutinum
\fi

\pars{Versus.}

\ifx\matversus\undefined
\ifx\matua\undefined
\else
\noindent \Vbardot{} Díriget Dóminus mansuétos in iudício. 

\noindent \Rbardot{} Docébit mites vias suas.
\fi
\ifx\matub\undefined
\else
\noindent \Vbardot{} Bonitátem et prudéntiam et sciéntiam doce me.

\noindent \Rbardot{} Quia præcéptis tuis crédidi.
\fi
\ifx\matuc\undefined
\else
\noindent \Vbardot{} Audiam quid loquátur Dóminus Deus.

\noindent \Rbardot{} Loquétur pacem ad plebem suam.
\fi
\ifx\matud\undefined
\else
\noindent \Vbardot{} Atténdite, pópule meus, doctrínam meam.

\noindent \Rbardot{} Inclináte aurem vestram in verba oris mei.
\fi
\else
\matversus
\fi

\ifx\sineabsolutio\undefined
\vspace{5mm}

\sineinitiali{temporalia/oratiodominica-mat.gtex}

\vspace{5mm}

\pars{Absolutio.}

\cuminitiali{}{temporalia/absolutio-ipsius.gtex}

\vfill
\pagebreak

\cuminitiali{}{temporalia/benedictio-solemn-deus.gtex}
\fi

\vspace{7mm}

\lectioi

\ifx\sineabsolutio\undefined
\noindent \Vbardot{} Tu autem, Dómine, miserére nobis.
\noindent \Rbardot{} Deo grátias.
\fi

\vfill
\pagebreak

\responsoriumi

\vfill
\pagebreak

\ifx\sineabsolutio\undefined
\cuminitiali{}{temporalia/benedictio-solemn-christus.gtex}
\fi

\vspace{7mm}

\lectioii

\ifx\sineabsolutio\undefined
\noindent \Vbardot{} Tu autem, Dómine, miserére nobis.
\noindent \Rbardot{} Deo grátias.
\fi

\vfill
\pagebreak

\responsoriumii

\vfill
\pagebreak

\ifx\sineabsolutio\undefined
\cuminitiali{}{temporalia/benedictio-solemn-ignem.gtex}
\fi

\vspace{7mm}

\lectioiii

\ifx\sineabsolutio\undefined
\noindent \Vbardot{} Tu autem, Dómine, miserére nobis.
\noindent \Rbardot{} Deo grátias.
\fi

\vfill
\pagebreak

\responsoriumiii

\vfill
\pagebreak

\rubrica{Reliqua omittuntur, nisi Laudes separandæ sint.}

\sineinitiali{temporalia/domineexaudi.gtex}

\vfill

\oratio

\vfill

\noindent \Vbardot{} Dómine, exáudi oratiónem meam.
\Rbardot{} Et clamor meus ad te véniat.

\vfill

\noindent \Vbardot{} Benedicámus Dómino.
\noindent \Rbardot{} Deo grátias.

\vfill

\noindent \Vbardot{} Fidélium ánimæ per misericórdiam Dei requiéscant in pace.
\Rbardot{} Amen.

\vfill
\pagebreak

\hora{Ad Laudes.} %%%%%%%%%%%%%%%%%%%%%%%%%%%%%%%%%%%%%%%%%%%%%%%%%%%%%

\cantusSineNeumas

\vspace{0.5cm}
\ifx\deusinadiutorium\undefined
\grechangedim{interwordspacetext}{0.18 cm plus 0.15 cm minus 0.05 cm}{scalable}%
\cuminitiali{}{temporalia/deusinadiutorium-communis.gtex}
\grechangedim{interwordspacetext}{0.22 cm plus 0.15 cm minus 0.05 cm}{scalable}%
\else
\deusinadiutorium
\fi

\vfill
\pagebreak

\ifx\hymnuslaudes\undefined
\ifx\hiemalislaudes\undefined
\ifx\lauda\undefined
\else
\pars{Hymnus}

\cuminitiali{I}{temporalia/hym-PergrataMundo.gtex}
\fi
\ifx\laudc\undefined
\else
\pars{Hymnus}

\cuminitiali{VIII}{temporalia/hym-PergrataMundo-einsiedeln.gtex}
\fi
\ifx\laudb\undefined
\else
\pars{Hymnus}

\cuminitiali{I}{temporalia/hym-AEterneLucis-hk.gtex}
\fi
\ifx\laudc\undefined
\else
\pars{Hymnus}

\cuminitiali{VIII}{temporalia/hym-PergrataMundo-einsiedeln.gtex}
\fi
\ifx\laudd\undefined
\else
\pars{Hymnus}

\cuminitiali{IV}{temporalia/hym-AEterneLucis.gtex}
\fi
\else
\ifx\laudac\undefined
\else
\pars{Hymnus}

\cuminitiali{I}{temporalia/hym-PergrataMundo.gtex}
\fi
\ifx\laudbd\undefined
\else
\pars{Hymnus}

\grechangedim{interwordspacetext}{0.16 cm plus 0.15 cm minus 0.05 cm}{scalable}%
\cuminitiali{IV}{temporalia/hym-AEterneLucis.gtex}
\grechangedim{interwordspacetext}{0.22 cm plus 0.15 cm minus 0.05 cm}{scalable}%
\vspace{-3mm}
\fi
\fi
\else
\hymnuslaudes
\fi

\vfill
\pagebreak

\ifx\laudes\undefined
\ifx\lauda\undefined
\else
\pars{Psalmus 1.} \scriptura{Ps. 23, 4.3}

\vspace{-4mm}

\antiphona{IV* e}{temporalia/ant-innocensmanibus.gtex}

\scriptura{Psalmus 23.}

\initiumpsalmi{temporalia/ps23-initium-iv_-e-auto.gtex}

\input{temporalia/ps23-iv_-e.tex} \Abardot{}

\vfill
\pagebreak

\pars{Psalmus 2.} \scriptura{Tob. 13, 6}

\vspace{-4mm}

\antiphona{VII a}{temporalia/ant-exaltateregem.gtex}

\scriptura{Canticum Tobiæ, Tob. 13, 2-8}

\initiumpsalmi{temporalia/tobiae-initium-vii-a-auto.gtex}

\input{temporalia/tobiae-vii-a.tex} \Abardot{}

\vfill
\pagebreak

\pars{Psalmus 3.} \scriptura{Ps. 32, 1; \textbf{H93}}

\vspace{-4mm}

\antiphona{IV E}{temporalia/ant-rectosdecet.gtex}

\vspace{-4mm}

\scriptura{Psalmus 32.}

\vspace{-2mm}

\initiumpsalmi{temporalia/ps32-initium-iv-E-auto.gtex}

\input{temporalia/ps32-iv-E.tex}

\vfill

\antiphona{}{temporalia/ant-rectosdecet.gtex}

\vfill
\pagebreak
\fi
\ifx\laudb\undefined
\else
\pars{Psalmus 1.} \scriptura{Ps. 42, 5; \textbf{H95}}

\vspace{-4mm}

\antiphona{VI F}{temporalia/ant-salutarevultusmei.gtex}

\scriptura{Psalmus 42.}

\initiumpsalmi{temporalia/ps42-initium-vi-F-auto.gtex}

\input{temporalia/ps42-vi-F.tex} \Abardot{}

\vfill
\pagebreak

\pars{Psalmus 2.} \scriptura{Is. 38, 20; \textbf{H95}}

\vspace{-7mm}

\antiphona{E}{temporalia/ant-cunctisdiebus.gtex}

\vspace{-4mm}

\scriptura{Canticum Ezechiæ, Is. 38, 10-20}

\vspace{-3mm}

\initiumpsalmi{temporalia/ezechiae-initium-e-auto.gtex}

\input{temporalia/ezechiae-e.tex} \Abardot{}

\vfill
\pagebreak

\pars{Psalmus 3.} \scriptura{Ps. 64, 2; \textbf{H96}}

\vspace{-4mm}

\antiphona{VIII a}{temporalia/ant-tedecet.gtex}

\vspace{-2mm}

\scriptura{Psalmus 64.}

\vspace{-2mm}

\initiumpsalmi{temporalia/ps64-initium-viii-A-auto.gtex}

\input{temporalia/ps64-viii-A.tex} \Abardot{}

\vfill
\pagebreak
\fi
\ifx\laudc\undefined
\else
\iffalse
\pars{Psalmus 1.} \scriptura{Ps. 83, 5}

\vspace{-4mm}

\antiphona{VIII G}{temporalia/ant-beatiquihabitant.gtex}

\vspace{-2mm}

\scriptura{Psalmus 84.}

\vspace{-2mm}

\initiumpsalmi{temporalia/ps84-initium-viii-G-auto.gtex}

\input{temporalia/ps84-viii-G.tex} \Abardot{}
\else
\pars{Psalmus 1.} \scriptura{Ps. 84, 2; \textbf{H97}}

\vspace{-4mm}

\antiphona{VI F}{temporalia/ant-benedixistidomine.gtex}

\vspace{-2mm}

\scriptura{Psalmus 84.}

\vspace{-2mm}

\initiumpsalmi{temporalia/ps84-initium-vi-F-auto.gtex}

\input{temporalia/ps84-vi-F.tex} \Abardot{}
\fi

\vfill
\pagebreak

\pars{Psalmus 2.}

\vspace{-4mm}

\antiphona{VII d}{temporalia/ant-denoctespiritusmeus.gtex}

\vspace{-2mm}

\scriptura{Canticum Isaiæ, Is. 26, 1-12}

\vspace{-2mm}

\initiumpsalmi{temporalia/isaiae3-initium-vii-d.gtex}

\input{temporalia/isaiae3-vii-d.tex} \Abardot{}

\vfill
\pagebreak

\pars{Psalmus 3.} \scriptura{Ps. 66, 2}

\vspace{-4mm}

\antiphona{E}{temporalia/ant-illuminadomine.gtex}

%\vspace{-2mm}

\scriptura{Psalmus 66.}

%\vspace{-2mm}

\initiumpsalmi{temporalia/ps66-initium-e.gtex}

\input{temporalia/ps66-e.tex} \Abardot{}

\vfill
\pagebreak
\fi
\ifx\laudd\undefined
\else
\pars{Psalmus 1.}

\vspace{-4mm}

\antiphona{VIII G}{temporalia/ant-cantaboetpsalmum.gtex}

\vspace{-2mm}

\scriptura{Psalmus 100.}

\vspace{-2mm}

\initiumpsalmi{temporalia/ps100-initium-viii-G-auto.gtex}

\input{temporalia/ps100-viii-G.tex} \Abardot{}

\vfill
\pagebreak

\pars{Psalmus 2.} \scriptura{Ps. 50, 19}

\vspace{-4mm}

\antiphona{I f}{temporalia/ant-sacrificiumdeo.gtex}

%\vspace{-2mm}

\scriptura{Canticum Danielis, Dan. 3, 26.27.29.34-41}

%\vspace{-2mm}

\initiumpsalmi{temporalia/dan32-initium-i-f-auto.gtex}

\input{temporalia/dan32-i-f.tex} \Abardot{}

\vfill
\pagebreak

\pars{Psalmus 3.} \scriptura{Ps. 143, 1; \textbf{H99}}

\vspace{-4mm}

\antiphona{VI F}{temporalia/ant-benedictusdominusdeusmeus.gtex}

%\vspace{-2mm}

\scriptura{Psalmus 143, 1-10.}

%\vspace{-2mm}

\initiumpsalmi{temporalia/ps143i_x-initium-vi-F-auto.gtex}

\input{temporalia/ps143i_x-vi-F.tex} \Abardot{}

\vfill
\pagebreak
\fi
\else
\laudes
\fi

\ifx\lectiobrevis\undefined
\ifx\lauda\undefined
\else
\pars{Lectio Brevis.} \scriptura{Rom. 13, 11.12-13}

\noindent Hora est iam vos de somno súrgere. Nox procéssit, dies autem appropiávit. Abiciámus ergo ópera tenebrárum et induámur arma lucis. Sicut in die honéste ambulémus.
\fi
\ifx\laudb\undefined
\else
\pars{Lectio Brevis.} \scriptura{1 Th. 5, 4-5}

\noindent Vos, fratres, non estis in ténebris, ut vos dies ille tamquam fur comprehéndat; omnes enim vos fílii lucis estis et fílii diéi. Non sumus noctis neque tenebrárum.
\fi
\ifx\laudc\undefined
\else
\pars{Lectio Brevis.} \scriptura{1 Io. 4, 14-15}

\noindent Nos vídimus et testificámur quóniam Pater misit Fílium salvatórem mundi. Quisque conféssus fúerit: Iesus est Fílius Dei, Deus in ipso manet, et ipse in Deo.
\fi
\ifx\laudd\undefined
\else
\pars{Lectio Brevis.} \scriptura{Is. 55, 1}

\noindent Omnes sitiéntes, veníte ad aquas; et, qui non habétis argéntum, properáte, émite et comédite, veníte, émite absque argénto et absque ulla commutatióne vinum et lac.
\fi
\else
\lectiobrevis
\fi

\vfill

\ifx\responsoriumbreve\undefined
\ifx\laudac\undefined
\else
\pars{Responsorium breve.} \scriptura{Ps. 17, 3} 

\cuminitiali{VI}{temporalia/resp-deusmeusadiutormeus.gtex}
\fi
\ifx\laudbd\undefined
\else
\pars{Responsorium breve.} \scriptura{Ps. 118, 149.147}

\cuminitiali{VI}{temporalia/resp-vocemmeamaudi.gtex}
\fi
\else
\responsoriumbreve
\fi

\vfill
\pagebreak

\ifx\benedictus\undefined
\ifx\laudac\undefined
\else
\pars{Canticum Zachariæ.} \scriptura{Lc. 1, 69; \textbf{H422}}

\vspace{-5mm}

{
\grechangedim{interwordspacetext}{0.18 cm plus 0.15 cm minus 0.05 cm}{scalable}%
\antiphona{VII a}{temporalia/ant-erexitnobisdominus.gtex}
\grechangedim{interwordspacetext}{0.22 cm plus 0.15 cm minus 0.05 cm}{scalable}%
}

%\vspace{-3mm}

\scriptura{Lc. 1, 68-79}

%\vspace{-1mm}

\cantusSineNeumas
\initiumpsalmi{temporalia/benedictus-initium-vii-a-auto.gtex}

\input{temporalia/benedictus-vii-a.tex} \Abardot{}
\fi
\ifx\laudbd\undefined
\else
\pars{Canticum Zachariæ.} \scriptura{Lc. 1, 71; \textbf{H423}}

\vspace{-5mm}

{
\grechangedim{interwordspacetext}{0.18 cm plus 0.15 cm minus 0.05 cm}{scalable}%
\antiphona{I g\textsuperscript{5}}{temporalia/ant-demanuomnium.gtex}
\grechangedim{interwordspacetext}{0.22 cm plus 0.15 cm minus 0.05 cm}{scalable}%
}

%\vspace{-3mm}

\scriptura{Lc. 1, 68-79}

%\vspace{-1mm}

\initiumpsalmi{temporalia/benedictus-initium-i-g5-auto.gtex}

\input{temporalia/benedictus-i-g5.tex} \Abardot{}
\fi
\else
\benedictus
\fi

\vspace{-1cm}

\vfill
\pagebreak

\pars{Preces.}

\sineinitiali{}{temporalia/tonusprecumnovum.gtex}

\ifx\preces\undefined
\ifx\lauda\undefined
\else
\noindent Vocatiónis cæléstis partícipes facti, fratres caríssimi, \gredagger{} Iesum, pontíficem confessiónis nostræ, \grestar{} benedicámus clamántes:

\Rbardot{} Dómine Deus noster et salvátor noster.

\noindent O Rex omnípotens, \gredagger{} qui, per baptísmum, regále nobis sacerdótium contulísti, \grestar{} fac ut laudis tibi semper sacrifícium offerámus.

\Rbardot{} Dómine Deus noster et salvátor noster.

\noindent Da nobis mandáta tua serváre, \grestar{} ut per Sanctum Spíritum in te maneámus et tu in nobis.

\Rbardot{} Dómine Deus noster et salvátor noster.

\noindent Sapiéntiam tuam da nobis ætérnam, \grestar{} ut nobíscum sit hódie et nobíscum operétur.

\Rbardot{} Dómine Deus noster et salvátor noster.

\noindent Concéde nobis hódie néminem umquam contristáre, \grestar{} omnes autem, qui nobíscum sunt, lætificáre.

\Rbardot{} Dómine Deus noster et salvátor noster.
\fi
\ifx\laudb\undefined
\else
\noindent Salvatóri nostro benedicámus, qui sua resurrectióne mundum clarificávit, \grestar{} et humíliter invocémus eum dicéntes:

\Rbardot{} Salva nos, Dómine, in sémita tua.

\noindent Resurrectiónem tuam, Dómine Iesu, oratióne cólimus matutína, \grestar{} spes glóriæ tuæ diem nostrum illúminet.

\Rbardot{} Salva nos, Dómine, in sémita tua.

\noindent Súscipe, Dómine, vota et propósita nostra, \grestar{} tamquam diéi nostri primítias.

\Rbardot{} Salva nos, Dómine, in sémita tua.

\noindent Tríbue in dilectióne tua nos hódie profícere, \grestar{} ut ómnia in nostrum omniúmque bonum cooperéntur.

\Rbardot{} Salva nos, Dómine, in sémita tua.

\noindent Da, Dómine, sic lucére lucem nostram coram homínibus, \grestar{} ut vídeant ópera nostra bona et Patrem gloríficent.

\Rbardot{} Salva nos, Dómine, in sémita tua.
\fi
\ifx\laudc\undefined
\else
\noindent Christum, qui sibi pópulum novum acquisívit in sánguine suo, \grestar{} adorémus et rogémus supplíciter:

\Rbardot{} Meménto pópuli tui, Dómine.

\noindent Rex et Redémptor noster, audi Ecclésiam tuam, huius diéi inítio te laudántem, \grestar{} doce eam sine intermissióne glorificáre maiestátem tuam.

\Rbardot{} Meménto pópuli tui, Dómine.

\noindent Spes et fírmitas nostra, in te confídimus, \grestar{} neque umquam confundémur.

\Rbardot{} Meménto pópuli tui, Dómine.

\noindent Infirmitátem nostram réspice et in adiutórium nostrum festína, \grestar{} quia sine te nihil póssumus fácere.

\Rbardot{} Meménto pópuli tui, Dómine.

\noindent Páuperum et derelictórum meménto, ne dies novus sit eis pondus, \grestar{} sed solámen et gáudium.

\Rbardot{} Meménto pópuli tui, Dómine.
\fi
\ifx\laudd\undefined
\else
\noindent Deus, lætítiam nobis concédens hoc matutíno témpore eum laudándi, spem róborat nostram. \grestar{} Ideo fidénter ipsi supplicémus:

\Rbardot{} Ad glóriam nóminis tui, audi nos, Dómine.

\noindent Grátias ágimus tibi, Deus et Pater Salvatóris nostri Iesu, \grestar{} pro cognitióne et immortalitáte, quas dedísti nobis per ipsum.

\Rbardot{} Ad glóriam nóminis tui, audi nos, Dómine.

\noindent Concéde humilitátem córdibus nostris, \grestar{} ut subiécti simus ínvicem in timóre Christi.

\Rbardot{} Ad glóriam nóminis tui, audi nos, Dómine.

\noindent Spíritum tuum nobis fámulis tuis infúnde, \grestar{} ut diléctio nostra fratérna sine simulatióne fiat.

\Rbardot{} Ad glóriam nóminis tui, audi nos, Dómine.

\noindent Qui hómini labórem præcepísti in mundi dominatiónem, \grestar{} concéde, ut labor noster te célebret fratrésque nostros sanctíficet.

\Rbardot{} Ad glóriam nóminis tui, audi nos, Dómine.
\fi
\else
\preces
\fi

\vfill

\pars{Oratio Dominica.}

\cuminitiali{}{temporalia/oratiodominicaalt.gtex}

\vfill
\pagebreak

\rubrica{vel:}

\pars{Supplicatio Litaniæ.}

\cuminitiali{}{temporalia/supplicatiolitaniae.gtex}

\vfill

\pars{Oratio Dominica.}

\cuminitiali{}{temporalia/oratiodominica.gtex}

\vfill
\pagebreak

% Oratio. %%%
\oratio

\vspace{-1mm}

\vfill

\rubrica{Hebdomadarius dicit Dominus vobiscum, vel, absente sacerdote vel diacono, sic concluditur:}

\vspace{2mm}

\antiphona{C}{temporalia/dominusnosbenedicat.gtex}

\rubrica{Postea cantatur a cantore:}

\vspace{2mm}

\ifx\benedicamuslaudes\undefined
\cuminitiali{IV}{temporalia/benedicamus-feria-laudes.gtex}
\else
\benedicamuslaudes
\fi

\vspace{1mm}

\vfill
\pagebreak

\end{document}

